\section{The Benchmark: Measuring a Live Merchant Ecosystem}
\label{sec:benchmark}

Most work on stopping problems assumes a price distribution. We would rather measure
one. This section describes the corpus, what it does and does not represent, and the
two properties that make it usable as a benchmark.

\subsection{Collection}

Merchants implementing the Universal Commerce Protocol publish a discovery document
and expose a read-only catalog search operation. We probed \ResultMerchantsTargeted{}
such merchants and paginated each catalog, collecting \ResultCatalogRows{} listings
covering \ResultUniqueSkus{} distinct SKUs. Of these, \ResultCrossMerchantSkus{} SKUs
appear at more than one independent merchant, which is what makes comparison possible
at all.

Only read operations were used; nothing was added to a cart and no purchase was
attempted. Merchant eligibility follows a documented permission ladder, and the
inventory used here is labelled \emph{declared capability}: the merchant advertises
the operation we exercised. That is a recorded modelling assumption, not a
case-by-case terms-of-service review, and it is a limitation rather than a warrant.

\subsection{The instrument does not add noise}

Before trusting any difference between two dates, we checked the scanner against
itself. Two full scans of the same merchants minutes apart returned byte-identical
listings and prices. Pagination is deterministic and result ordering is stable, so
any difference observed between snapshots is drift in the world.

This check matters more than it might appear. An early comparison suggested roughly
half of all listings vanished within a month; it was an artifact of comparing a
shallow scan against a deep one. Restricting the comparison to merchants whose
catalogs were paginated \emph{in full} on both dates removes the artifact, and we
report only that restricted figure.

\subsection{Offers move}

Across \ResultTrackedListings{} listings tracked on fully paginated merchants,
\ResultDelistingRate{} were delisted between the two dates and
\ResultPriceChangeRate{} of the survivors changed price, with a median relative move
of \ResultMedianRelativePriceChange{} among those that moved. Prices are sticky but
not static, and disappearance is the more common event. An agent that defers a
purchase is therefore taking a real risk, which is the premise the rest of the paper
depends on.

\subsection{Episodes}

An episode is one exact SKU offered by at least two independent merchants in a single
currency. The earlier snapshot supplies merchant price and availability forecasts;
the later snapshot is frozen as the held-out panel that every policy is replayed
against. Because the split is chronological, no policy can see its own evaluation
data.

Two filters do most of the work and both cost us episodes. Cross-currency listings
are dropped, since comparing a euro price to a dollar price requires an exchange-rate
assumption we prefer not to make; this removes the majority of candidate SKUs. A SKU
missing from the later snapshot counts as delisted only when that merchant was
paginated in full, and otherwise the merchant is dropped from the episode rather than
recorded as a stockout. What survives is \ResultEpisodeCount{} episodes, of which
\ResultHeldOutEpisodeCount{} are held out for evaluation.

\subsection{What this corpus is not}

The resulting benchmark is narrow in a way worth stating plainly. The median episode
has \ResultMedianMerchantCount{} merchants and a price ratio of
\ResultMedianPriceDispersionRatio{}, meaning that in half of all episodes the
merchants charge nearly the same price and no policy can distinguish itself. Thin
cross-merchant overlap appears to be a property of this ecosystem rather than of our
pipeline: even in the deeper snapshot, almost every shared SKU is carried by exactly
two merchants. This is a finding about the ecosystem, and it directly bounds how much
any stopping rule can achieve here.
